
Like \texttt{Flatten}, \texttt{PreFlatten} runs on the \texttt{SSA} IR:
currently, it runs before \texttt{Flatten}, so that the lattice-based
\texttt{Flatten} analysis has a better chance to pick up on opportunities to
completely flatten particular uses.

\section{Motivation}
The basic intuition is that if we see a sequence like:
\begin{minted}{python}
  def f(...) -> ...:
    ...other blocks...

    block bbN():
      t: 'a * 'b = tuple(x0, x1)
      # call g(t) and return to some other block bbM
      call bbM g(t)

    ...other blocks...

  def g(t: 'a * 'b) -> ...:
    block bb0():
      x0: 'a = select(t, 0)
      x1: 'b = select(t, 1)
      ...etc...

    ...other blocks...
\end{minted}
then we can eliminate the intermediate tuple \texttt{t} with a simple peephole
optimization that ``pushes down'' the tuple allocation from the caller into the
callee, i.e.:
\begin{minted}{python}
  def f(...) -> ...:
    ...other blocks...

    block bbN():
      # call g(x0, x1) and return to some other block bbM
      call bbM g_flat(x0, x1)

    ...other blocks...

  def g_flat(x0_flat: 'a, x1_flat: 'b) -> ...:
    block bb0():
      t: 'a * 'b = tuple(x0_flat, x1_flat)
      x0: 'a = select(t, 0)
      x1: 'b = select(t, 1)
      ...etc...

    ...other blocks...
\end{minted}
and then rely on existing optimization machinery to \hl{remove the allocation
  completely}.

In case all callers of \texttt{g} satisfy the pattern above, then the result of
this transformation is equivalent to the existing \texttt{Flatten} pass.
However, since we want to phrase this optimization as a simple, purely-local
transformation, we need a solution for the case when not all callers of
\texttt{g} are \texttt{PreFlatten}-able: we resolve this problem by emitting a
fresh definition \texttt{g\_flat} and leaving the existing \texttt{g} in place
for callers that require it, i.e. for something like:
\begin{minted}{python}
  def f(t1: 'a * 'b) -> ...:
    ...other blocks...

    block bbN1():
      # call g(t1) and return to some other block bbM
      call bbM g(t1)

    block bbN2():
      t2: 'a * 'b = tuple (x1, x2)
      # call g(t2) and return to some other block bbM
      call bbM g(t2)

    ...other blocks...

 def g(t: 'a * 'b) -> ...:
    ...as above...
\end{minted}
We can emit:
\begin{minted}{python}
  def f(t1: 'a * 'b) -> ...:
    ...other blocks...

    block bbN1():
      # Unflattenable caller: no tuple allocation to ``push down''
      call bbM g(t1)

    block bbN2():
      # Flattenable caller: can ``push down'' the allocation
      call bbM g_flat(x1, x2)

    ...other blocks...

 def g_flat(x0_flat: 'a, x1_flat: 'b) -> ...:
    ...as above...

 def g(t: 'a * 'b) -> ...:
    ...as above...
\end{minted}
By iteratively applying this strategy, we can ``push together'' distant pairs of
flatten/unflatten operations, and so try to ensure that the maximum possible
number of callsites in the program are flattened -- even in cases where the
global \texttt{Flatten} transformation would fail.

This strategy also naturally extends to eliminating \texttt{ConApp}s: starting
from a pattern like:
\begin{minted}{python}
  datatype Env = Env1 of 'a * 'b | ...
  def f(...) -> ...:
    ...
    block bbN():
      t: 'a * 'b = ...
      env: Env = con Env1 (t)
      call bbM g(env)
    ...

  def g(e: Env) -> ...:
    block bb0():
      case e of Env1 => goto bb1 | ...
\end{minted}
we can specialize the callsite and ``push down'' the \texttt{ConApp} to arrive at:
\begin{minted}{python}
  def f(...) -> ...:
    ...
    block bbN():
      t: 'a * 'b = ...
      call bbM g(t)
    ...

  def g_flat_env1(t: 'a * 'b) -> ...:
    block bb0():
      e: Env = con Env1 (t)
      case e of Env1 => goto bb1 | ...
\end{minted}
and, again, rely on existing optimization patterns to eliminate the
\texttt{ConApp} completely, in addition to elimating now-dead blocks
corresponding to \texttt{Env} cases that are no longer reachable, etc. For
higher order functions, this subsumes the action of the existing
\texttt{Polyvariance} pass.

\hl{TODO: CONAPP IS NOT INTRODUCED PROPERLY BEFORE THIS}

\section{Implementation}
The two type flattening implementations share a common framework, differing only
in the types of operations that they match on:

\begin{enumerate}
\item Find all \texttt{call} statements with a ``pack'' operation that feeds
  directly into a \texttt{call}: for the \texttt{tuple} case, this corresponds
  to \texttt{tuple} constructor applications:
\begin{minted}{python}
  t: a * b * ... = tuple(x1, x2, ...)
  call bbN f_tuple(t, y, z, ...) # call1
  ....
  def f_tuple(t: a * b * ..., y: c, z: d, ...) -> ...:
\end{minted}
and for the \texttt{datatype} case, this corresponds to \texttt{ConApp}s:
\begin{minted}{python}
  datatype ConType = Con1 of a * b * ... | Con2 of ...
  ...
  t: a * b * ... = ...
  ...
  c: ConType = con Con1(t)
  call bbN f_con(c, y, z, ...) # call2
  ...
  def f_con(c: ConType, y: c, z: d, ...) -> ...:
\end{minted}
\item Tag each such \texttt{call} statement with a per-argument flattening
  decision, to indicate which arguments are valid targets to flatten, and
  associate the tagged \texttt{call}s with their callee. For the flattenable
  arguments, we include the information necessary to construct the flattened
  call (i.e. the arguments to the \texttt{tuple} / \texttt{ConApp} call).
  Continuing the examples from above, we would have:
\begin{minted}{python}
    f_tuple_caller_info = {
      'call1': [FlattenVar(x1, x1, ...), PreserveVar, PreserveVar, ...]
      ...other callers...
    }
    f_con_caller_info = {
      'call2': [FlattenVar(t), PreserveVar, PreserveVar, ...]
      ...other callers...
    }
\end{minted}
\item Traverse all function-callee-decision collections and apply heuristics to
  ``downgrade'' \texttt{FlattenVar} decisions to \texttt{PreserveVar} decisions
  when specialization would be unprofitable.

\item Traverse all function-callee-decision collections and emit the new
  required specialized versions: an all-\texttt{PreserveVar} decision
  corresponds to the original function. Any \texttt{FlattenVar(...)} decisions
  result in a new specialization whose formal parameters include the flat
  arguments in the appropriate location, and whose first block re-binds the
  flattened names to the old, un-flattened name before jumping to the original
  first block. That is, for the \texttt{tuple} example above, we emit:
\begin{minted}{python}
  # Specialization for [FlattenVar(...), PreserveVar, ...]
  def f_tuple_flat0(x1: a, x2: b, ..., y: c, z: d, ...) -> ...:
    block bbBind():
      # Re-introduce the original, unflattened name
      t: a * b * ... = tuple(x1, x2, ...)
      # Jump to the original first block
      goto bb0()
\end{minted}
and for the \texttt{ConApp} example, we emit:
\begin{minted}{python}
  # Specialization for [FlattenVar(...), PreserveVar, ...]
  def f_con_flat0(t: a * b * ..., ..., y: c, z: d, ...) -> ...:
    block bbBind():
      # Re-introduce the original, unflattened name
      c: ConType = con Con1(t)
      # Jump to the original first block
      goto bb0()
\end{minted}

\item Traverse the program to visit all \texttt{call} statements again, and
  redirect the call target to the appropriate specialization that we identified
  in the previous step (perhaps the original function). For non-identity
  flattening decisions, use the contents of the \texttt{FlattenVar(...)}
  argument to construct the appropriate flat call, i.e. for \texttt{tuple}:
\begin{minted}{python}
  t: a * b * ... = tuple(x1, x2, ...)
  call bbN f_tuple(t, y, z, ...) # call1
\end{minted}
becomes:
\begin{minted}{python}
  call bbN f_tuple_flat0(x1, x2, ..., y, z, ...)
\end{minted}
and for \texttt{ConApp}, 
\begin{minted}{python}
  c: ConType = con Con1(t)
  call bbN f_con(c, y, z, ...) # call2
\end{minted}
becomes:
\begin{minted}{python}
  call bbN f_con_flat0(t, y, z, ...)
\end{minted}

\item With the callsite-flattening mapping still intact and with the duplicated
  functions retaining the original variable names, iteratively re-apply the
  process above on all newly-emitted functions until either convergence or a
  user-specified iteration limit is reached. This fixes an issue with (mutually)
  recursive functions, where a collection of $n$ mutually recursive calls
  require $n$ iterations to remove unnecessary pack/unpack operations, e.g. for
  a source program like:
\begin{minted}{python}
  f(...) -> ...:
    ...
    # flat call to `g`
    t0: a * b = (x, y)
    call bbN g(t0)

  g(t: a * b) -> ...:
    ...
    block bbM():
      # arguments for recursive call
      t2: a * b = (w, z)
      call bbK g(t2)
\end{minted}
The action of the pass is as follows: for both the call in \texttt{f} and
\texttt{g}, we can flatten the function argument, and so we emit a new
specialization of the function, \texttt{g\_flat}
\begin{minted}{python}
  f(...) -> ...:
    ...
    # flat call to `g`
    t0: a * b = (x, y)
    call bbN g_flat(x, y)  # pointed to flat specialization

  g(t: a * b) -> ...:
    ...
    block bbM():
      call bbK g_flat(w, z) # pointed to flat specialization

  g_flat(x: a, y: b) -> ...:
    block bbBind():
      # reintroduce original, unflattened name
      t: a * b = tuple(x, y)
      # jump to original first block
      goto bb0()

    ...
    block bbM():
      # arguments for recursive call
      t2: a * b = (w, z)
      # This call is not flattened, but it could be!
      call bbK g(t2)
\end{minted}
In this step, we re-run the flattening logic using the already-emitted set of
specializations to remove the unflattened call in the specialized function, i.e.:
\begin{minted}{python}
  g_flat(x: a, y: b) -> ...:
    block bbBind():
      # reintroduce original, unflattened name
      t: a * b = tuple(x, y)
      # jump to original first block
      goto bb0()

    ...
    block bbM():
      call bbK g_flat(w, z) # pointed to flat specialization
\end{minted}
Multiple iterations are required to handle the case of mutually recursive
functions, in a similar way. \hl{TODO: EXAMPLE?}

\item Introduce fresh names for the symbols (specifically block labels and
  variable names) in all duplicated functions. This is a requirement from MLton,
  which expects that symbols' names are globally unique.

\item Signal either ``progress'' (if we emitted a new specialization) or
  ``done'' to the caller. Repeat until either convergence or a user-specified
  iteration limit is reached.
\end{enumerate}

\section{Cost Heuristics}

\hl{TODO: SHOULD RUN EXPERIMENTS WHERE WE RUN AFTER FLATTEN AND ALSO ONLY
  FLATTEN NON-TAIL CALLS}

%% \subsection{Tuple-flattening}

%% \subsection{Constructor-flattening}

\section{Limitations}
Note that the worst-case bloat in IR size is quadratic in a single iteration:
the very upper bound is that every statement in the program causes us to
duplicate the entire program. This is clearly problematic in practice, and it
interacts particularly poorly with chunkification (since flattened tuples can be
more expensive than non-flat tuples to pass across chunk boundaries, due to the
required stack writes that are linear in the number of tuple elements).

\hl{TODO: EXPAND}

